\section{Conclusions and Future Work}
\label{sec:conclusions}

This project investigated whether metadata-aware retrieval can improve
Text-to-Spark-SQL generation over a heterogeneous data lake.

The implemented prototype combines a structured metadata catalog, local
semantic embeddings, relation-aware metadata retrieval, SQL-safe context
rendering, a local language model, SQL validation, optional one-shot repair,
and Apache Spark SQL execution.

The central design choice is to avoid providing the complete metadata catalog
to the language model for every request. Instead, the system retrieves and
expands only the metadata considered relevant to the current analytical
question, while explicitly preserving physical datasets, columns,
relationships, and semantic rules required for execution.

The experimental results support this approach.

On the frozen held-out benchmark, complete-catalog prompting achieved 70\%
result accuracy. Relation-aware metadata retrieval increased result accuracy
to 80\%, and the addition of validation-driven one-shot repair increased it
further to 85\%.

At the same time, relation-aware prompting reduced mean prompt size from
2558.45 to 776.35 tokens, corresponding to a 69.66\% reduction. In the
controlled CPU-only latency experiment, mean prompt-evaluation time decreased
from 182.58 seconds to 54.24 seconds, a reduction of 70.29\%.

The catalog-scale experiment showed that the retrieved context remains
approximately bounded while the complete metadata catalog grows. The full
catalog increased from 9,498 characters with four sources to 16,651
characters with sixteen sources, while the mean retrieved context remained
close to 2,100 characters.

This scalability benefit is accompanied by an important limitation. Macro
metadata recall decreased from 0.925 to 0.825 after semantically related
hard-negative sources were introduced. Error analysis showed that the
degradation was caused primarily by competition around geographic and weather
metadata rather than by catalog size alone.

The held-out results also demonstrated the importance of retrieval
completeness. Queries with perfect metadata recall achieved substantially
higher downstream result accuracy than queries with missing required
metadata. Validation and repair can correct some invalid SQL generations, but
they cannot reliably compensate for relationships or schema information that
were never retrieved.

The Spark data-volume experiment evaluated the physical execution layer from
100,000 rows up to the complete 2,963,713-row curated Yellow Taxi dataset.
The temporal taxi--weather workload showed the clearest runtime growth, while
the simpler workloads were strongly influenced by fixed execution and warm-up
costs in the two-core local environment.

Overall, the results suggest that metadata management is a central component
of data-centric generative-AI systems for structured analytics. Improving the
way metadata is selected and represented can improve both correctness and
efficiency without changing the underlying language model.

Future work could extend the prototype in several directions.

First, retrieval could use an adaptive Top-\(k\) strategy or confidence-aware
expansion rather than a fixed dense retrieval depth. This could improve recall
when semantically similar catalog sources compete for the same retrieval
budget.

Second, structural reranking could make stronger use of catalog relationships,
schema compatibility, cardinality information, and expected analytical
operations before the final context is rendered.

Third, the metadata catalog could incorporate richer provenance, data-quality
statistics, source freshness, and additional semantic constraints.

Fourth, the evaluation could be expanded to larger and more diverse
cross-domain benchmarks, including catalogs that contain substantially more
physical sources and unrelated business domains.

Fifth, the Spark execution experiment could be repeated on a distributed
cluster and with much larger physical datasets in order to study partitioning,
shuffle behaviour, executor scaling, and resource utilisation.

Finally, larger or alternative local language models could be evaluated using
the same frozen metadata and benchmark infrastructure. This would make it
possible to separate improvements caused by stronger language models from
those caused by better metadata grounding.

The main contribution of the project is therefore not a claim that
relation-aware retrieval eliminates Text-to-SQL errors. Instead, the project
provides a reproducible demonstration that metadata selection, structural
expansion, SQL-safe representation, validation, and controlled evaluation can
substantially influence the behaviour of language-model-based analytics over
heterogeneous Big Data sources.
